\section{Robustness}\label{sec:robustness}

I vary one part of the model at a time: the preparation-cost law, the value classes and the prior, and the
challenger's information. Appendix~\ref{app:robustness} states each result in full and proves
it.

\subsection{Preparation costs}\label{sec:costs}

The rise of entry in Proposition~\ref{prop:ccc} survives other preparation-cost laws, with or without atoms
(Proposition~\ref{prop:atomless}). The cost now mixes a low component with weight \(\rho\) and a high
component in place of the two-point law of Section~\ref{sec:values}; either component may be atomless or have
atoms. The proposition keeps parts (i) and (ii) of Proposition~\ref{prop:ccc} under the trading-cost window (H2)
when every cost in the low component lies strictly below \(B_{F_1}(m(\pi))\) and every cost in the high component
lies strictly between \(B_{F_0}(\pi)\) and \(B_{F_1}(M(\pi))\). The low component then prepares at every feasible
belief, so entry stays at least \(\rho\), and the trading bounds of Section~\ref{sec:informed} apply unchanged.

At the benchmark, the low component runs from zero to \(\reg{atom_L}\), and the high component is uniform on
\([c_H-\varepsilon,c_H+\varepsilon]\). Every half-width below \(\reg{p4_atomless_max_width_floor_r2}\) keeps all
three parts of Proposition~\ref{prop:ccc}. With half-width \(\varepsilon=\reg{atom_eps}\),
entry at \(r_1\) is \(\reg{p4_atomless_r1_entry}\), the value under the two-point law. At \(r_2\), gross profit at
the top posterior, which lies \(\reg{margin_return_B}\) below \(c_H\), falls inside this band, and entry at \(r_2\)
is \(\reg{p4_atomless_r2_entry}\). With an atomless high component, the drop at the entry ceiling becomes a ramp:
full-order entry is continuous and nonincreasing in strength.

\subsection{Value classes and priors}\label{sec:classes}

The reversal holds for every pair of ranked classes and every prior. The classes satisfy the class
ranking~\eqref{eq:classes} and have different means, and the reserve, the prior, the zero-cost probability and the
slot are arbitrary. For uniform incumbents, a nonempty open set of strengths and costs then satisfies (H1)--(H3)
and the ownership condition of Proposition~\ref{prop:ccc} (Proposition~\ref{prop:classes}).

The construction puts the weak and the strong incumbent
just above the lowest value \(v_*\) at which the classes start to differ above the reserve. The spread vanishes near
\(v_*\), so informed trade against the weaker incumbent cannot cover its cost. The profit gap stays near the full class
gap there, so a high cost that only a favorable price justifies fits inside it. In the point-mass case
\(v_*=\ell\), so the reversal holds for every \(h>\ell\).

Table~\ref{tab:classes} in Appendix~\ref{app:robustness} reports six economies with their own costs and
strengths. With wide uniform classes, entry rises from \(\reg{p2_ex_c_r0_entry}\) to
\(\reg{p2_ex_c_r1_entry}\), and high-class ownership rises although the high class can lose to the strong
incumbent; point masses at the same means leave entry at \(\reg{p2_ex_d_r1_entry}\). A prepared challenger keeps
\((\theta-\max\{p,R\})_+\), a convex function of its value, so a mean-preserving spread within a class weakly
raises gross profit at every belief.

With more than two classes, the reversal holds under a separation condition (Proposition~\ref{prop:manyclasses}).
The price reveals only the posterior mean of class proceeds. Separation asks for a gap between the
proceeds of each high-group class and every feasible posterior mean, and the same gap for the low group. In the
strong equilibrium, class proportions within each group then stay at their prior proportions, so the price reveals
the whole posterior vector.

A four-class economy near the benchmark satisfies separation with distance \(\reg{p2_groups_gamma}\).
Without separation, one price can carry two posterior vectors: a strict three-class
equilibrium has two orders with different posterior vectors and one price.

\subsection{Complementary private signals}\label{sec:signals}

The challenger can gain from the price even when its own information is more accurate than the investor's. A buyer
may know its integration plans, while investors who follow the target know its customers or its product market.
Proposition~\ref{prop:signals} represents the two sources with conditionally independent binary
signals. The investor's signal has accuracy \(a_T\), and the challenger's signal has accuracy \(d\).

The zero-cost challenger still prepares after every price and signal, so the price still reveals the public
posterior, and the challenger combines that posterior with its own signal. Under strict conditions on the costs,
the weak incumbent has zero orders and entry \(\rho\); the strong incumbent has full orders and entry above
\(\rho\). The conditions allow \(d>a_T\).

The signal example uses its own parameter vector \(S(B)\). It keeps the slot, the values, the reserve and the
prior of the benchmark, and it sets \(\rho=\reg{sig_rho}\), \(c_H=\reg{sig_cH}\), \(k=\reg{sig_k}\),
\(r_0=\reg{sig_r0}\) and \(r_1=\reg{sig_r1}\). With investor accuracy \(a_T=\reg{sig_aT}\) and challenger accuracy
\(d=\reg{sig_d}\), entry rises from \(\reg{p4_s_signal_r0_entry}\) to \(\reg{p4_s_signal_r1_entry}\).
A favorable signal alone does not justify high-cost preparation against the weak incumbent;
a favorable signal and a favorable price together do so against the strong incumbent. Both sources are needed:
without the challenger's signal, entry is \(\reg{p4_s_no_y_r0_entry}\) at both strengths.

With the accuracies of the signal example \(S(B)\), the benchmark has a unique equilibrium at each strength when
an indifferent investor may mix over its best orders. At \(r_0\), the investor mixes between
trade and abstention after either signal, and it trades with probability \(\reg{p4_b_signal_r0_buy_rate}\). At
\(r_1\), it buys with probability one after a favorable signal, and after an unfavorable signal it sells with
probability \(\reg{p4_b_signal_r1_sell_rate}\). Entry falls from \(\reg{p4_b_signal_r0_entry}\) to
\(\reg{p4_b_signal_r1_entry}\). When an investor that is indifferent between abstention and a nonzero order must
abstain, neither strength has an equilibrium.

At the benchmark itself, the conditions hold with \((a_T,d)=(\reg{sig_comp_aT},\reg{sig_comp_d})\).
Entry then rises from \(\reg{p4_comp_b_r0_entry}\) to \(\reg{p4_comp_b_r1_entry}\), and
without the challenger's signal it is \(\reg{p4_comp_b_no_y_r1_entry}\) at both strengths. A buyer signal at least
as accurate as the investor's needs other costs and strengths, as in the signal example: at the benchmark the
strong trading bound needs an accurate investor signal, and the weak bound then fails.
